\section{Building Blocks}
\label{sec:blocks}

SNARK constructions typically reduce circuit satisfiability, expressed through constraint systems such as R1CS or PLONK, to matrix or vector relations involving inner products, Hadamard products, and permutations.
These relations can be encoded as polynomial identities and proved using sumcheck, or proved using other specialized protocols, such as folding-based inner-product arguments~\cite{sp/BunzBBPWM18}.
In this section, we assume prover protocols for inner and Hadamard products over \(\ExtensionField\) whose prover costs are linear in the vector length, measured in extension-field operations.
We develop a unified vector polynomial packing technique that reduces base-field relations to shorter extension-field relations, yielding bitwidth awareness by applying these protocols to the packed instances.
In Section~\ref{sec:together}, we combine these building blocks into a SNARK protocol for PLONK constraints.

\subsection{Packed Inner-Product Proof}
\label{sec:ip}

We introduce vector polynomial packing by the packed inner-product proof that reduces a
length-\(N\) inner-product claim over \(\Field\) to a
length-\(m\) claim over \(\ExtensionField\), where \(N=m\kappa\).
Consider the claim:
\[
    c=\langle\bm a,\bm b\rangle,
    ~
    \bm a,\bm b\in\Field^N.
\]
We partition each vector into \(\kappa\) blocks of length \(m\), and
arrange these blocks as:
\begin{equation}
\label{eq:matAB}
    \bm A
    \coloneqq
    \concat{
        \bm a^{(1)},\ldots,\bm a^{(\kappa)}
    },
    \qquad
    \bm B
    \coloneqq
    \concat{
        \bm b^{(1)},\ldots,\bm b^{(\kappa)}
    }
    \in\Field^{m\times\kappa}.
\end{equation}
We use the monomial basis
\(
    \bm{\phi}_A(X)
    =
    (1,X,\ldots,X^{\kappa-1})
\)
and the reciprocal monomial basis
\(
    \bm{\phi}_B(X)
    =
    (1,X^{-1},\ldots,X^{-(\kappa-1)})
\)
to pack the columns into a polynomial vector and a Laurent
polynomial vector, respectively:
\[
    \bm g_A(X)
    \coloneqq
    \sum_{j=1}^{\kappa}\bm a^{(j)}X^{j-1}
    \in\Field[X]^m,
    \qquad
    \bm g_B(X)
    \coloneqq
    \sum_{j=1}^{\kappa}\bm b^{(j)}X^{-(j-1)}
    \in\Field[X^{-1}]^m.
\]
Their inner product defines the Laurent polynomial:
\begin{equation}
    f(X)
    \coloneqq
    \left\langle
        \bm g_A(X),\bm g_B(X)
    \right\rangle.
\label{eq:ip-laurent-polynomial}
\end{equation}
Its constant coefficient is
\(\sum_{j=1}^{\kappa}
\langle\bm a^{(j)},\bm b^{(j)}\rangle
=\langle\bm a,\bm b\rangle\),
while all cross-column terms have nonzero exponents.
Thus, \(c=\langle\bm a,\bm b\rangle\) if and only if there
exist \(h_+,h_-\in\Field[X]\), each of degree at most
\(\kappa-2\), such that:
\begin{equation}
    f(X)
    =
    c
    +Xh_+(X)
    +X^{-1}h_-(X^{-1}).
\label{eq:ip-correction-identity}
\end{equation}

With \(c\) and the commitments to \(\bm A\) and \(\bm B\)
fixed, the prover sends the coefficients of \(h_+\) and \(h_-\).
The verifier then samples
\(\ext{x}\sample\ExtensionField^\times\).
The prover evaluates both encodings at \(\ext{x}\) to obtain:
\[
    \ext{\bm a}
    \coloneqq
    \bm g_A(\ext{x}) \in \ExtensionField^m,
    \qquad
    \ext{\bm b}
    \coloneqq
    \bm g_B(\ext{x})
    \in\ExtensionField^m.
\]
Evaluating Eq.~\eqref{eq:ip-correction-identity} at \(\ext{x}\)
reduces the polynomial identity to the length-\(m\)
inner-product claim:
\begin{equation}
    \left\langle
        \ext{\bm a},\ext{\bm b}
    \right\rangle
    =
    c
    +\ext{x}h_+(\ext{x})
    +\ext{x}^{-1}h_-(\ext{x}^{-1}).
\label{eq:packed-inner-product-check}
\end{equation}
The verifier computes the RHS from \(c\),
the auxiliary polynomials, and \(\ext{x}\).
The parties then apply the extension-field inner-product
protocol of Fact~\ref{fact:ip}, reducing this claim to two
MLE claims on the length-\(m\) packed vectors.
These terminal claims are authenticated against the
original witness commitments via a packed PCS.
Algorithm~\ref{alg:ip} in Appendix~\ref{appx:pseudocode}
gives the packed inner-product protocol.

\begin{proposition}[Packed inner-product proof]
\label{proposition:ip-complexity}
Algorithm~\ref{alg:ip} is a proof protocol for the base-field inner-product relation
with soundness error at most:
\[
    \frac{2\kappa-2}{|\ExtensionField|-1}
    +
    \frac{2\log_2 m}{|\ExtensionField|}.
\]
Excluding the terminal MLE openings, its prover time cost is
\(
    O(N \kappa + \kappa^3)\bb
\)
and its peak prover memory is \(O(N+\kappa^2) \Field\).
\end{proposition}

We prove Proposition~\ref{proposition:ip-complexity}
in Appendix~\ref{appx:proof-ip}.


\subsection{Packed Hadamard-Product Proof}
\label{sec:hadamard}

We next apply vector polynomial packing, using a different choice of
basis polynomials, to a weighted
Hadamard-product relation:
\begin{equation}
    \ext{c}
    =
    \left\langle
        \ext{\bm\xi},
        \bm a\circ\bm b
    \right\rangle \in \ExtensionField,
    \qquad
    \ext{\bm\xi}
    =
    \ext{\bm r}\kronecker\ext{\bm l} \in \ExtensionField^N,
\label{eq:hadamard-claim}
\end{equation}
where
\(
    \bm a,\bm b\in\Field^N,~
    \ext{\bm l}\in\ExtensionField^m,~
    \ext{\bm r}\in\ExtensionField^\kappa,
\)
and
\(
    N=m\kappa.
\)
When expressed using MLE polynomials, this weighted
Hadamard relation corresponds to the degree-two
zerocheck setting considered by Dao and Thaler~\cite{iacr/DaoT24}
and Gruen~\cite{iacr/Gruen24}.

The tensor-structured weight vector factors as
\(\ext{\bm\xi}=\ext{\bm r}\kronecker\ext{\bm l}\).
Using the block matrices
\(\bm A,\bm B\in\Field^{m\times\kappa}\)
defined in Eq.~\eqref{eq:matAB}, we can therefore rewrite
Eq.~\eqref{eq:hadamard-claim} in matrix form as:
\[
    \ext{c}
    =
    \ext{\bm l}^{\top}
    (\bm A\circ\bm B)
    \ext{\bm r}
    \in\ExtensionField.
\]

To prove this identity, the prover sends
\(
    \ext{\bm c}^{\top}
    \coloneqq
    \ext{\bm l}^{\top}(\bm A\circ\bm B) \in \ExtensionField^\kappa.
\)
The verifier checks
\(\ext{c}\EqOrNot
\langle\ext{\bm c},\ext{\bm r}\rangle\).
We then encode these \(\kappa\) column claims at distinct
interpolation nodes to verify them jointly,
where the $\kappa$ entries of $\ext{c}$ is denoted by $\ext{c}_1, \cdots, \ext{c}_\kappa \in \ExtensionField$.

We choose both \(\bm{\phi}_A(X)\) and \(\bm{\phi}_B(X)\)
to be the Lagrange basis
\((\psi_1(X),\ldots,\psi_\kappa(X))\)
from Definition~\ref{def:lagrange}, and let \(\Psi(X)\)
denote the corresponding vanishing polynomial.
We define:
\[
\begin{aligned}
    \bm g_A(X)
    &\coloneqq
    \sum_{j=1}^{\kappa}
    \bm a^{(j)}\psi_j(X),
    \quad
    \bm g_B(X)
    \coloneqq
    \sum_{j=1}^{\kappa}
    \bm b^{(j)}\psi_j(X),
    \quad
    g_C(X)
    \coloneqq
    \sum_{j=1}^{\kappa}
    \ext{c}_j\psi_j(X).
\end{aligned}
\]
When \(|\Field|\ge 2\kappa-1\), we choose \(\kappa\) distinct
interpolation nodes in \(\Field\), so the Lagrange basis
polynomials have coefficients in \(\Field\).
Consequently:
\[
    \bm g_A(X),\bm g_B(X)\in\Field[X]^m,
    \qquad
    g_C(X)\in\ExtensionField[X].
\]

We define the residual polynomial:
\[
    f(X)
    \coloneqq
    \left\langle
        \ext{\bm l},
        \bm g_A(X)\circ\bm g_B(X)
    \right\rangle
    -
    g_C(X)
    \in\ExtensionField[X].
\]
At each interpolation root \(\gamma_j\in\mathcal H\),
the property of the Lagrange basis gives:
\[
    f(\gamma_j)
    =
    \left\langle
        \ext{\bm l},
        \bm a^{(j)}\circ\bm b^{(j)}
    \right\rangle
    -
    \ext{c}_j.
\]
Thus, all \(\kappa\) column claims hold if and only if
\(f\) vanishes on \(\mathcal H\), or equivalently,
\(\Psi(X)\) divides \(f(X)\).
Since \(\deg f\le 2\kappa-2\) and \(\deg\Psi=\kappa\),
this is equivalent to the existence of
\(h(X)\in\ExtensionField[X]\) with \(\deg h\le\kappa-2\)
such that
\begin{equation}
    f(X)=\Psi(X)h(X).
\label{eq:hadamard-quotient}
\end{equation}

For valid claims, the values of \(f\) on \(\mathcal H\)
are known to be zero.
When \(|\Field|\ge 2\kappa-1\), the prover can therefore
recover \(f\) by evaluating it at \(\kappa-1\) distinct
points in \(\Field\setminus\mathcal H\) and interpolating
these values together with the known zeros.
Dividing the resulting polynomial by \(\Psi\) then gives \(h\).

After the prover sends \(\ext{\bm c}\) and \(h\), the verifier
samples
\(
    \ext{x}
    \sample
    \ExtensionField\setminus\mathcal H.
\)
The Lagrange evaluations give the packing weights:
\[
    \ext{\bm\rho}(\ext{x})
    \coloneqq
    \bigl(
        \psi_1(\ext{x}),\ldots,\psi_\kappa(\ext{x})
    \bigr)^\top \in \ExtensionField^\kappa.
\]
The prover computes the packed vectors:
\[
\begin{aligned}
    \ext{\bm a}
    &\coloneqq
    \bm A\ext{\bm\rho}(\ext{x})
    =
    \bm g_A(\ext{x}) \in \ExtensionField^m,
    \qquad
    \ext{\bm b}
    \coloneqq
    \bm B\ext{\bm\rho}(\ext{x})
    =
    \bm g_B(\ext{x})  \in \ExtensionField^m,
\end{aligned}
\]
and the verifier directly computes a degree-$(\kappa-1)$ polynomial value:
\[
    \ext{c}_x
    \coloneqq
    g_C(\ext{x})
    =
    \left\langle
        \ext{\bm c},
        \ext{\bm\rho}(\ext{x})
    \right\rangle \in \ExtensionField.
\]
Evaluating Eq.~\eqref{eq:hadamard-quotient} at \(\ext{x}\)
reduces the Hadamard-product claim to:
\begin{equation*}
    \left\langle
        \ext{\bm a},
        \ext{\bm b}\circ\ext{\bm l}
    \right\rangle
    =
    \ext{c}_x
    +
    \Psi(\ext{x})h(\ext{x}).
\end{equation*}
This is again a length-\(m\) inner-product relation over
\(\ExtensionField\).
Applying the inner-product protocol of Fact~\ref{fact:ip}
reduces it to two MLE claims on
\(\ext{\bm a}\) and \(\ext{\bm b}\).

Algorithm~\ref{alg:hadamard} in Appendix~\ref{appx:pseudocode}
specifies the protocol.

\iffalse
Fact~\ref{fact:ip} reduces this length-\(m\) inner product to two
MLE claims. Write its challenge point and terminal values as:
\[
    \ext{\bm y}\in\ExtensionField^{\log_2m},
    \qquad
    \ext{v}_1,\ext{v}_2\in\ExtensionField.
\]
The terminal claims on the original vectors are:
\begin{align}
\label{eq:hadamard-terminal}
    \ext{v}_1
    &\EqOrNot
    \left\langle
        \bm a,
        \ext{\bm\rho}(\ext{x})
        \kronecker
        \tensorspan{\ext{\bm y}}
    \right\rangle,
    \qquad
    \ext{v}_2
    \EqOrNot
    \left\langle
        \bm b,
        \ext{\bm\rho}(\ext{x})
        \kronecker
        \left(
            \ext{\bm l}
            \circ
            \tensorspan{\ext{\bm y}}
        \right)
    \right\rangle.
\end{align}

The second row weight is represented by multiplying the corresponding
two-entry factors, as in Eq.~\eqref{eq:tensor-weight-product};
it need not be a normalized MLE weight.
Both claims fit the packed-opening interface in
Eq.~\eqref{eq:pcs-opening-relation}.
Algorithm~\ref{alg:hadamard} in Appendix~\ref{appx:pseudocode}
additionally aggregates the column claims at a fresh MLE point,
retaining these row factors.
\fi

\smallskip\noindent\textbf{Small fields via an intermediate field.}
When \(q=|\Field|<\kappa\), we choose the interpolation roots
in an intermediate field.
Starting from the desired extension degree \(\kappa\),
choose \(t=2\lceil\log_q\kappa\rceil-1\), increasing \(t\)
if necessary until
\(
    q^t\ge  2 t\left\lceil\frac{\kappa}{t}\right\rceil - 1.
\)
Then round \(\kappa\) up to the nearest multiple of \(t\).
We reuse \(\kappa\) for the adjusted degree, which satisfies
\(t\mid\kappa\), \(q^t\ge 2\kappa-1\), and
\(t=O(\log_q\kappa)\).
This gives the field tower:
\[
    \Field\subseteq\mathbb K=\mathbb F_{q^t}
    \subseteq\ExtensionField,
    \qquad
    [\ExtensionField:\mathbb K]=\kappa/t.
\]

Choose \(\mathcal H\subseteq\mathbb K\) of size \(\kappa\)
and define \(\bm\psi(X),\Psi(X),\bm g_A(X),\bm g_B(X),g_C(X)\) as above.
After zero padding if needed, arrange the witnesses as
\(\bm A,\bm B\in\Field^{m\times\kappa}\), with \(N=m\kappa\).
The witness entries remain over \(\Field\), while:
\[
    \bm g_A (X),\bm g_B (X)\in\mathbb K[X]^m,
    \qquad
    g_C(X)\in\ExtensionField[X].
\]

We implement the interpolation procedure using
precomputed addition and multiplication tables over
\(\mathbb K\), which also handle inputs in
\(\Field\times\mathbb K\).
Together, these tables require \(O(t^2q^{2t})\)
base-field preprocessing operations and \(O(tq^{2t})\)
base-field coordinates of storage.
Caching the evaluation and interpolation data requires
an additional \(O(\kappa^2t)\) coordinates.
We charge each table lookup as $O(1)$.

Under this convention, evaluating the witness encodings
and aggregating their row products costs \(O(N\kappa)\).
Evaluating \(g_C\), interpolating \(f\), and computing
the quotient cost \(O(\kappa^3/t)\).
Including the remaining protocol steps, but excluding PCS
commitments and openings, the online prover time and
total storage are
\[
    T_{\Prover}=O(N\kappa+\kappa^3),
    \qquad
    S_{\Prover}=O(N+\kappa^2t+tq^{2t}),
\]
measured in base-field operations and coordinates.
The storage bound includes all precomputed data and is
achieved by processing sample points and witness rows
sequentially.
For \(m\ge\kappa\), the online prover time simplifies
to \(O(N\kappa)\).

The packed polynomial identity, challenge sampling over
\(\ExtensionField\), and terminal extension-field inner-product
proof carry over to the adjusted parameters.
The same soundness and complexity bounds for the
remaining steps therefore apply.

\begin{proposition}[Packed Hadamard-product proof]
\label{proposition:hadamard-complexity}
Let \(q=|\Field|\), \(N=m\kappa\), and
\(t=[\mathbb K:\Field]\), where \(\mathbb K\) is the
interpolation field and may equal \(\Field\).
Algorithm~\ref{alg:hadamard} is complete for the weighted
Hadamard-product relation.
With the terminal MLE claims verified, its soundness error
is at most:
\[
    \frac{2\kappa-2}{|\ExtensionField|-\kappa}
    +
    \frac{2\log_2N}{|\ExtensionField|}.
\]
With the precomputed data described above, its online
prover cost, excluding terminal MLE openings, is
\(
    O(N\kappa+\kappa^3)\,\bb.
\)
Its peak prover memory, including all precomputed data, is
\(
    O\bigl(N+\kappa^2t+(t-1)q^{2t}\bigr) \Field.
\)
\end{proposition}

We prove Proposition~\ref{proposition:hadamard-complexity} in
Appendix~\ref{appx:proof-hadamard}.



\subsection{Packed Permutation Proof}
\label{sec:permutation}

We next consider the permutation relation:
\[
    \bm a\stackrel{\mathsf{perm}}{=}\bm b,
    \qquad
    \bm a,\bm b\in\Field^N.
\]
It asserts that the two vectors contain the same multiset of
entries. Such proofs verify the PLONK copy constraints in
Eq.~\eqref{eq:plonk2} and are used in lookup
constructions~\cite{eurocrypt/SettyTW24}.
After the witness commitments are fixed, the verifier samples
\(\ext z\sample\ExtensionField\). The prover claims a common
fingerprint \(\ext\Theta\) and proves:
\begin{equation}
    \ext\Theta
    =\mathsf{GrandProd}(\ext z\bm1+\bm a)
    =\mathsf{GrandProd}(\ext z\bm1+\bm b).
\label{eq:packed-permutation-fingerprint}
\end{equation}
The two fingerprint polynomials are identical if and only if
the permutation relation holds. It therefore suffices to prove
the two grand-product claims. We describe the proof for \(\bm a\).

\smallskip\noindent\textbf{Grand products from Hadamard products.}
Arrange \(\bm a\) as \(\bm A\in\Field^{m\times\kappa}\), as in
Eq.~\eqref{eq:matAB}, with \(N=m\kappa\).
Let \(q=|\Field|\), and assume that \(m,\kappa\) are powers
of two with \(m\geq\kappa\geq2\).
We first construct a polynomial product tree over \(\Field\),
keeping \(\kappa\) columns and halving the number of rows at each
level:
\begin{equation}
\begin{aligned}
    \bm C_0(Z)&\coloneqq Z \cdot \bm1_{m\times\kappa}+\bm A,\\
    \bm C_{\tau+1}(Z)
      &\coloneqq\bm C_{\tau,U}(Z)\circ\bm C_{\tau,D}(Z),
      \qquad \tau \in [0,\log_2\kappa-1].
\end{aligned}
\label{eq:packed-permutation-product-tree}
\end{equation}
where \(\bm C_{\tau,U}(Z)\) and \(\bm C_{\tau,D}(Z)\) are the upper
and lower row halves of \(\bm C_\tau(Z)\). Thus:
\[
    \bm C_\tau(Z)\in
       \Field[Z]^{(m/2^\tau)\times\kappa},
    \qquad
    \deg_Z \bm C_\tau(Z)=2^\tau.
\]
The prover computes these polynomial coefficients over \(\Field\), then
evaluates \(\bm C_{\log_2\kappa}(Z)\) using precomputed powers of
\(\ext z\) and $\Field-\ExtensionField$ products.
Its \(m\) entries form the vector
\(\bm c_{\log_2\kappa}(\ext z)\in\ExtensionField^m\).
The remaining levels are computed over
\(\ExtensionField\):
\[
    \bm c_{\tau+1}(\ext z)
       =\bm c_{\tau,U}(\ext z)\circ\bm c_{\tau,D}(\ext z),
    \qquad
    \tau \in [\log_2\kappa, \log_2 N-1].
\]
where \(U,D\) now denote vector halves. The final entry is
\(\bm c_{\log_2N}(\ext z)=\ext\Theta\).

The proof traverses these trees in reverse. Starting from
\(\ext\Theta\), weighted Hadamard proofs over \(\ExtensionField\)
reduce the last \(\log_2m\) levels to an MLE claim on
\(\bm c_{\log_2\kappa}(\ext z)\).
Factoring its tensor weight into column and row weights gives:
\[
    \ext v_{\log_2\kappa}
      =\ext{\bm l}_{\log_2\kappa}^{\top}
         \bm C_{\log_2\kappa}(\ext z)
         \ext{\bm r}_{\log_2\kappa}.
\]
At each subsequent level, the recursive claim has the same form:
\[
    \ext v_{\tau+1}
       =\ext{\bm l}_{\tau+1}^{\top}
          \bm C_{\tau+1}(\ext z)\ext{\bm r}_{\tau+1},
    \qquad
    \ext{\bm l}_{\tau+1}\in\ExtensionField^{m/2^{\tau+1}},
    \quad
    \ext{\bm r}_{\tau+1}\in\ExtensionField^\kappa,
\]
with a tensor-structured row weight and fixed column weights.

\smallskip\noindent\textbf{Base-field interpolation: \(q\geq2\kappa-1\).}
In this case, the interpolation points lie in \(\Field\).
We use the Lagrange basis from Section~\ref{sec:hadamard}:
\[
    \bm\psi(X)\coloneqq
       (\psi_1(X),\ldots,\psi_\kappa(X))^\top,
    \qquad
    \bm g_\tau(Z,X)\coloneqq\bm C_\tau(Z)\bm\psi(X).
\]
Choose \(\kappa-1\) distinct points
\(\gamma'_1,\ldots,\gamma'_{\kappa-1}\in
\Field\setminus\mathcal H\).
For each level \(\tau\in [\log_2\kappa]\) and point
\(\gamma'_j\), fix \(X=\gamma'_j\) and retain \(Z\) as a
variable. Define the polynomial vector:
\[
\begin{aligned}
    \bm\eta_{\tau,j}(Z)
      &\coloneqq
      \bm g_{\tau-1,U}(Z,\gamma'_j)
      \circ\bm g_{\tau-1,D}(Z,\gamma'_j)
      =\sum_{\nu=0}^{2^\tau}\bm p_{\tau,j,\nu}Z^\nu,
    \quad
    \bm p_{\tau,j,\nu}\in\Field^{m/2^\tau}.
\end{aligned}
\]
Each entry has degree at most \(2^\tau\) in \(Z\).
In particular, since \(\sum_j\psi_j(X)=1\):
\[
    \bm g_0(Z,\gamma'_j)
       =Z \cdot \bm1_m+\bm A\bm\psi(\gamma'_j).
\]
The prover computes the coefficients
\(\bm p_{\tau,j,\nu}\) over \(\Field\). The weighted
polynomial is:
\[
    F_{\tau,j}(Z)
       \coloneqq\left\langle\ext{\bm l}_\tau,
                             \bm\eta_{\tau,j}(Z)\right\rangle
       =\sum_{\nu=0}^{2^\tau}
          \left\langle\ext{\bm l}_\tau,
                       \bm p_{\tau,j,\nu}\right\rangle Z^\nu
       \in\ExtensionField[Z].
\]
The prover first aggregates the coefficient vectors
\(\bm p_{\tau,j,\nu}\) against \(\ext{\bm l}_\tau\), obtaining
the extension-field coefficients of \(F_{\tau,j}(Z)\).
It then evaluates this scalar polynomial at \(\ext z\).
The powers \(1,\ext z,\ldots,\ext z^\kappa\) are shared across
all levels.

For the transition from \(\bm C_{\tau+1}\) to \(\bm C_\tau\),
where \(\tau \in [0,\log_2\kappa-1]\), the prover applies
Section~\ref{sec:hadamard} by sending the column claims:
\[
    \ext{\bm u}_{\tau+1}^{\top}
       \coloneqq
    \ext{\bm l}_{\tau+1}^{\top}
        \bigl(\bm C_{\tau,U}(\ext z)
               \circ\bm C_{\tau,D}(\ext z)\bigr).
\]
For each column, the prover first aggregates the coefficients
of \(\bm C_{\tau+1}(Z)\) against \(\ext{\bm l}_{\tau+1}\),
then evaluates the resulting scalar polynomial at \(\ext z\).
The verifier checks
\(\ext v_{\tau+1}=\ext{\bm u}_{\tau+1}^{\top}
\ext{\bm r}_{\tau+1}\).
As in Section~\ref{sec:hadamard}, define:
\[
\begin{aligned}
    g_C(X)&\coloneqq\ext{\bm u}_{\tau+1}^{\top}\bm\psi(X),\\
    f(X)&\coloneqq
       \left\langle\ext{\bm l}_{\tau+1},
          \bm g_{\tau,U}(\ext z,X)
             \circ\bm g_{\tau,D}(\ext z,X)\right\rangle
       -g_C(X).
\end{aligned}
\]
For valid column claims, \(f(X)\) vanishes on \(\mathcal H\).
Its off-domain samples are:
\[
    f(\gamma'_j)
       =F_{\tau+1,j}(\ext z)-g_C(\gamma'_j).
\]
Interpolating them together with the zeros on \(\mathcal H\)
gives \(f(X)\), and division by \(\Psi(X)\)) gives
\(h(X)\in\ExtensionField[X]\) with \(\deg h (X)\leq\kappa-2\).
The prover fixes \(\ext{\bm u}_{\tau+1}\) and \(h(X)\) before the
verifier samples a fresh
\(\ext x\sample\ExtensionField\setminus\mathcal H\).
The resulting check is:
\[
    \left\langle\ext{\bm l}_{\tau+1},
       \bm g_{\tau,U}(\ext z,\ext x)
         \circ\bm g_{\tau,D}(\ext z,\ext x)\right\rangle
       =g_C(\ext x)+\Psi(\ext x)h(\ext x).
\]
This is a weighted Hadamard claim of length \(m/2^{\tau+1}\)
over \(\ExtensionField\).

\smallskip\noindent\textbf{Recursive queries.}
We implement each extension-field weighted Hadamard proof by the degree-2 zero-check via
sumcheck over $\ExtensionField$. The checked polynomial has individual degree at most
three, and the two child claims use the same row MLE point
\(\ext{\bm y}\).
After their values \(\ext v_{\tau,U},\ext v_{\tau,D}\) are fixed,
the verifier samples a fresh \(\ext y_0\sample\ExtensionField\).
The next claim is obtained by setting:
\[
\begin{aligned}
    \ext v_\tau
      &\coloneqq(1-\ext y_0)\ext v_{\tau,U}
                    +\ext y_0\ext v_{\tau,D},
                    \qquad
    \ext{\bm l}_\tau
      \coloneqq\tensorspan{\ext y_0\mathbin\|\ext{\bm y}},
    \qquad
    \ext{\bm r}_\tau
      \coloneqq\ext{\bm\rho}(\ext x).
\end{aligned}
\]
It has the desired form
\(\ext v_\tau=\ext{\bm l}_\tau^\top
\bm C_\tau(\ext z)\ext{\bm r}_\tau\).
The same packed Hadamard-product proof is applied.
At level zero, both weight vectors sum to one, leaving the
claim:
\[
    \ext v_0-\ext z
       =\ext{\bm l}_0^\top\bm A\ext{\bm r}_0.
\]

\smallskip\noindent\textbf{Cost over \(\Field\).}
At level \(\tau\), forming the \(\kappa-1\) encodings from
\(\bm C_{\tau-1}\) costs \(O(N\kappa)\,\bb\).
Their row products cost \(O(m\kappa2^\tau)\,\bb\), totaling
\(O(N\kappa)\,\bb\) across all levels.
Coefficient aggregation costs \(O(N)\,\be\) per level,
including the column claims. The aggregated polynomials have
degree at most \(2^\tau\), so evaluating them at \(\ext z\)
costs \(O(\kappa2^\tau)\,\ee\) per level and
\(O(\kappa^2)\,\ee\) in total.
To form the terminal Hadamard vectors, the prover first
contracts the coefficients of \(\bm C_\tau(Z)\) with
\(\ext{\bm\rho}(\ext x)\), then evaluates the resulting row
polynomials at \(\ext z\), at a cost of
\(O(N)\,\be+O(m)\,\ee\) per level.
Using shared powers of \(\ext x\), interpolation and public
query generation cost \(O(\kappa^3)\,\bb\) per level.
The short Hadamard proofs and extension-field tree have
total length \(O(m)\).
Thus, under schoolbook arithmetic, the prover cost for one
grand-product proof is
\(
    O(N\kappa\log_2\kappa+\kappa^4)\bb.
\)


\smallskip\noindent\textbf{Intermediate-field tables over smaller fields: \(q<2\kappa-1\).}
Choose \(t\) and $\kappa$ such that $q^t \ge 2\kappa-1$ and $t \mid \kappa$, and we can use the tower:
\[
    \Field\subseteq\mathbb K=\mathbb F_{q^t}
       \subseteq\ExtensionField,
    \qquad [\ExtensionField:\mathbb K]=\kappa/t.
\]
We choose \(\mathcal H\) and the additional interpolation nodes
in \(\mathbb K\). The coefficients of \(\bm C_\tau(X)\) remain
in \(\Field\), while those of \(\bm g_\tau(X)\) and
\(\bm\eta_{\tau,j}(X)\) lie in \(\mathbb K\).
The polynomial identities and recursive queries are unchanged.

We precompute a multiplication table over \(\mathbb K\)
and count each table lookup as an \(O(1)\)-time operation.
Representing an extension-field element by \(\kappa/t\)
coordinates over \(\mathbb K\), a mixed
\(\mathbb K\times\ExtensionField\) product costs \(O(\kappa)\).
Using this multiplication table and precomputed Lagrange basis
evaluations, the grand-product reduction over \(\mathbb K\)
similarly achieves an online prover cost of
\(
    O\!\left(N\kappa\log_2\kappa+\frac{\kappa^4}{t}\right)\,\bb.
\)

Algorithm~\ref{alg:permutation} in Appendix~\ref{appx:pseudocode}
specifies the protocol for both cases.

\begin{proposition}[Packed permutation proof]
\label{proposition:permutation-complexity}
Let \(q=|\Field|\) be fixed and let
\(N=m\kappa\), where \(m,\kappa\) are powers of two and
\(m\geq\kappa\geq2\).
Algorithm~\ref{alg:permutation} is a proof protocol for the grand-product relation
with the soundness error:
\[
\begin{aligned}
    \epsilon_{\mathsf{GP}}\leq{}&
    \frac{\frac32\log_2m(\log_2m-1)+\log_2m}
         {|\ExtensionField|}+\sum_{\tau=0}^{\log_2\kappa-1}
      \left(
        \frac{2\kappa-2}{|\ExtensionField|-\kappa}
        +\frac{3\log_2(m/2^{\tau+1})+1}{|\ExtensionField|}
      \right).
\end{aligned}
\]
Consequently, when both input vectors are committed, the
permutation argument has soundness error at most:
\[
    \epsilon_{\mathsf{Perm}}\leq{} \min\left\{1,\frac{N-1}{|\ExtensionField|}\right\}
       +2\epsilon_{\mathsf{GP}}+\epsilon_{\mathsf{PCS}},
\]
where \(\epsilon_{\mathsf{PCS}}\) is the joint soundness error
of authenticating the two terminal packed-opening claims.
With the intermediate-field multiplication table and block
The online prover cost for one
grand product is
\(
    O\!\left(N\kappa\log_2\kappa+\kappa^4\right)\,\bb.
\)
The peak prover memory, including all precomputed tables, is:
\(
    O(N+\kappa^3+tq^{2t})\Field
\)
\end{proposition}

Proof of this proposition is provided in Appendix~\ref{appx:proof-permutation}.

\subsection{Packed Polynomial Commitment Scheme}
\label{sec:pcs}

We complete the packed reductions by authenticating their
terminal claims against a fixed-basis commitment.
Let \(N=m\kappa\) and \(\bm A\in\Field^{m\times\kappa}\).
The preceding protocols produce claims of the form:
\[
    \ext v
    =\langle\bm A\ext{\bm\rho},\ext{\bm\xi}\rangle
    =\ext{\bm\xi}^{\top}\bm A\ext{\bm\rho},
    \qquad
    \ext{\bm\rho}\in\ExtensionField^\kappa,
    \quad
    \ext{\bm\xi}\in\ExtensionField^m,
\]
where \(\ext{\bm\xi}\) is a tensor-structured row weight.
The query weights \(\ext{\bm\rho}\) may differ from the fixed
\(\Field\)-basis \(\ext{\bm\beta}\in\ExtensionField^\kappa\)
used to commit to the packed vector:
\[
    \ext{\bm a}=\bm A\ext{\bm\beta}\in\ExtensionField^m.
\]

To relate the two packings, the prover sends the column evaluations
\(\ext{\bm u}^{\top}\coloneqq\ext{\bm\xi}^{\top}\bm A\).
The verifier checks
\(\ext v\EqOrNot\langle\ext{\bm u},\ext{\bm\rho}\rangle\).
Write the query and the column claims in base-field coordinates:
\[
    \ext{\bm\xi}=\bm\Xi\ext{\bm\beta},
    \qquad
    \ext{\bm u}^{\top}=\ext{\bm\beta}^{\top}\bm U,
    \qquad
    \bm\Xi\in\Field^{m\times\kappa},
    \quad
    \bm U\in\Field^{\kappa\times\kappa}.
\]
Since \(\bm A\) has base-field entries, the column claims are
equivalent to \(\bm U=\bm\Xi^\top\bm A\).
Packing this matrix identity with the fixed basis gives:
\[
    \bm U\ext{\bm\beta}
    =\bm\Xi^\top\bm A\ext{\bm\beta}
    =\bm\Xi^\top\ext{\bm a}\in\ExtensionField^\kappa.
\]
After \(\ext{\bm u}\) is fixed, a fresh random projection reduces
this vector identity to a length-\(m\) inner-product claim involving
the committed vector \(\ext{\bm a}\).
Fact~\ref{fact:ip} then leaves one MLE opening of
\(\ext{\bm a}\) and one public query evaluation.
The verifier computes the latter directly from the tensor factors
of \(\ext{\bm\xi}\).

This is the ring-switching technique~\cite{eurocrypt/DiamondP26}
expressed in our vector and matrix notation.
Appendix~\ref{appx:pcs} gives the construction and its analysis;
Algorithm~\ref{alg:packed-pcs} specifies the protocol.
Under schoolbook arithmetic and \(m\geq\kappa\), the packed
opening reduction uses \(O(\kappa N)\,\bb\) and
\(O(N)\Field\) memory,
given an underlying extension-field PCS that commits and opens in \(O(m)\ee\) and uses \(O(m)\ExtensionField\) of memory.
